\begin{frame}{커널 함수: 내적을 우회한다}
  가장 널리 쓰이는 커널:
  \begin{columns}[T]
    \begin{column}{0.5\textwidth}
      \kb{다항 커널}(polynomial)
      \[
      K(x,x') = (x\cdot x' + c)^d
      \]
      ``도수 $d$ 이하의 \textbf{모든 단항식}''을 내포.
      $(x\cdot x'+1)^d$ 를 전개하면 도수 $d$ 이하의 모든
      단항식 $x_1^{a_1}x_2^{a_2}\cdots$ ($a_1+a_2+\cdots\le d$)
      이 (일정한 계수로) 전부 섞여 나온다.
    \end{column}
    \begin{column}{0.5\textwidth}
      \kb{RBF(가우시안) 커널}
      \[
      K(x,x') = \exp\Bigl(-\frac{\|x-x'\|^2}{2\sigma^2}\Bigr)
      \]
      두 점이 가까울수록 1에, 멀수록 0에 가까운 \textbf{유사도}
      함수. \textbf{무한 차원}의 $\phi$ 에 대응한다는 것이
      알려져 있는데도, 커널 함수 자체는 \textbf{유클리드 거리
      하나로} 계산.
    \end{column}
  \end{columns}
  \vskip0.4em
  {\footnotesize $\exp(-\gamma\|x-x'\|^2)$ 로 정의하는 문헌도
  많다 — $\gamma = 1/(2\sigma^2)$ 로 대응. scikit-learn의
  \texttt{SVC}는 \texttt{gamma}를 받는다.}
\end{frame}
